\chapter{Discussion}
\label{ch:discussion}

This chapter interprets the results of Chapter~\ref{ch:results}; it does not repeat them.
The thread running through the chapter is that the extractor loses transitions in two places, at
the nodes and at the edges, and that each level limits what the other can achieve.

\section{Key findings}
\label{sec:findings}

\paragraph{Two places where extraction fails.}
Static extraction of a transition graph from Angular source fails in two places, and the first
turned out to be the tool's rather than the method's.
Node extraction is complete on all six applications against the readers of the route
configuration (Table~\ref{tab:coverage}).
The runtime route oracle observes no node placed at an address where the framework mounts nothing.
That shows that no such error was observed, not that none can occur.
The runtime route oracle also bounds the first claim.
The route configuration it reads holds \ProbeMissed{} mounted addresses with no node,
\ProbeMissedGhostfolio{} of them in ghostfolio (Section~\ref{sec:res-probe}).
So only \ProbeExactApps{} of the six are complete in the stronger sense of a node for every mounted
address.
Both readings are reported.
The weaker is what the loss classification is computed against.
The stronger is what a user of the graph would want.
The node defects (constant paths, lazy boundaries) are traced in Section~\ref{sec:res-defects}.
The final version reads their causes (Section~\ref{sec:res-round2}), and those losses are empty.

\paragraph{What remains is lost at the edge.}
The navigation sites exist in the source and the screens they would connect exist in the graph,
but the extractor cannot decide which screens a site belongs to.
The exceptions are sites that the diagnosis labels as having a target screen with no node.
The label also covers targets held in a field.
Most of these sites are in ghostfolio (Table~\ref{tab:remaining}).
Scope attribution addresses this loss at the edge.
Under the order in which the loss rules are applied it is the largest loss category in every application
(Table~\ref{tab:loss}).
That order, however, is part of the measure (Section~\ref{sec:limitations}).
The ranking is a fact about an extractor, not a limit of the approach.
The two levels are not independent.
Attribution cannot draw an edge into a screen that has no node.
A complete set of nodes yields few transitions while the sites outside the component a route renders
stay unattributed.
Reporting only one level would have described a tool that works better than it does, which is why
Chapter~\ref{ch:results} reports node coverage beside must conversion.

\paragraph{RQ3 and what the loss rules can delimit.}
A lost construct is assigned to the first loss category that applies (Table~\ref{tab:loss-rules}).
A construct inside a shared component that appears only in screens without nodes is therefore
counted as an attribution loss, although attribution alone cannot recover it.
The categories are a diagnostic partition of the losses, not a bound on what attribution recovers.

The demotion rules decide whether the edges attribution added at k8sdash (its toolbar edges,
Section~\ref{sec:res-soundness}) occur on every input.
\CensusMustKdash{} of them are kept at must.
The demotion rules were written during development, in part in response to objections from red-team panels run with the AI assistant.
The witness check finds \MustWitnessViolations{} failures among the \MustWitnessEdges{} must edges.
It cannot read the destination of \MustWitnessUnread{}, all of them k8sdash's toolbar edges.
So only about half of that application's must edges are read.

\paragraph{RQ4 is conditional, and the condition is RQ1.}
Leave aside the rename operator, whose detection is close to guaranteed by construction.
The other three operators are each generated or scored only at constructs the graph already
represents in the simplest form the extractor reads.
They are detected in \MutFallibleDetected{} of \MutFallibleScored{} scored mutants on the attributed
graph: every deleted \texttt{routerLink}, \MutRetargetDetected{} of \MutRetargetScored{} retargeted
calls and \MutGuardDetected{} of \MutGuardScored{} removed guards (Table~\ref{tab:mutation}).
The mutation harness was corrected during development after its misses had been read
(Appendix~\ref{app:repairs}), so the clean result is not an independent test.
It does not show that the remaining operators would survive a larger or differently drawn sample.

For guards, what separates a detected mutant from an unobservable one is whether any edge enters a
screen the guarded route yields, which is the exclusion rule (Section~\ref{sec:res-mutation}).
The class ``guards unreliable in both directions'' of Chapter~\ref{ch:results} rests instead on
what Section~\ref{sec:res-round2} measured.
These are a guard carried by a route a wrapper function returns, and a \texttt{canActivate} element
written as a call.
Source review in red-team panels found those defects, not the witness check.
The mutation run neither confirms nor refutes the class.

Detection is also conditional on recovery.
The conclusion is therefore as follows.
Given that a construct is represented in the graph, every scored change to it was visible.
The probability that it is represented at all is what RQ1 measures.
For change impact analysis this form of claim is appropriate, because an impact query is only asked
about the part of the application that the model represents.
The graph states which part through its witnesses.

\paragraph{Two uses of the graph.}
The first, with measured support, is change impact analysis.
Every edge carries the file and line of the construct that produced it.
Changed lines therefore map to the edges whose witnesses fall inside them.
From there they map to the screens those edges leave and enter.
This is the screen-level counterpart of Chianti \cite{ren2004chianti}, in the impact analysis
tradition of Arnold and Bohner \cite{arnold1993impact,bohner1996software,lehnert2011taxonomy}.
RQ4 shows that the mapping reacts to the injected edits for the constructs the graph represents.
That is a property such an analysis needs, not a measure of its accuracy, which would ask whether
every affected screen is named and how many are named in excess.

The second, not itself measured, is checking in continuous integration at no execution cost.
The graph is recomputed from source on each commit, deterministically and without installing, serving or crawling anything (Section~\ref{sec:protocol}).
Properties such as ``a screen still has an inbound deterministic transition'' can then be asserted like
a type check. A guard assertion is the member of that list the graph cannot yet carry, and for a screen no edge enters there is nothing to read.
A dynamically built model cannot be used this way at the same cost, because it needs a running
application, test data and an exploration budget \cite{crawljax2012,feedex2013,websurvey2025}.

\paragraph{Where both uses stop.}
Both stop at runtime dynamic routing: a target computed from a value, a route table assembled at
startup from configuration, and an overlay no route describes.
There the static graph is silent or over-approximates into the may tier, which carries no soundness
claim.
Under the assumptions of Claim~\ref{claim:must} the must tier is a lower bound on the navigation an
application can perform; the may tier is not.
A query that must not miss a transition therefore cannot be answered from the graph alone.
The graph's witnesses and declared approximations tell a consumer which queries those are.
Authentication bounds the deterministic tier in the same way.
A guard decides on state the source does not fix, so a reader that keeps Claim~\ref{claim:must}
demotes it.
The small must tier is in part this fact and not a defect of the extractor.
Taking every guard to let the navigation through raises the tier from \CensusMust{} to
\GuardOnlyMust{} edges (Section~\ref{sec:res-round2}).

\section{Limitations and threats to validity}
\label{sec:limitations}

The results are measurements of one tool, on one corpus, with five methods and an end-to-end check
on two applications, performed by the person who wrote the tool.
Each of those four choices limits what the figures support.

\textbf{The tool was developed together with this evaluation, on the same corpus.} uigraph was built
for this thesis and corrected while it was evaluated (Appendix~\ref{app:repairs}). Every figure measured on the six applications is therefore in-sample.
The extractor, the independent route parser, the mutation harness and the structural reference were
corrected against the applications they are scored on. The structural reference guided two repairs of the extractor. The end-to-end check comes closest to an out-of-sample test. The extractor was fixed at its final commit before the suites' logs were read, and nothing was
repaired after.
The archive cannot corroborate that order, since the preregistration and the logs were committed
together. The corpus figures should therefore not be read as estimates for other applications.

\textbf{The scope is compared in both directions on two applications only, and only against the analysis's own relations.}
Every other method in Chapter~\ref{ch:evaluation} checks a node, an edge or a witness.
None can be asked whether the set of screens attribution computes for a site is the set it fires
from.
Three methods bear on it, each partly.
\begin{itemize}[nosep]
\item The hand-read structural reference (Section~\ref{sec:res-groundtruth}) answers in both
directions, but for two applications of six and only as agreement with the analysis's own
relations. On alfresco the full analysis recalls \GTAlfFullRecallNonShell{} of the pairs off the shell
(\GTAlfFullRecall{} overall) at \GTAlfFullPrecision{} precision, below the class join's
\GTAlfClassPrecision{} (Section~\ref{sec:res-groundtruth}). On realworld it finds \GTFullRecall{} of its (site, screen) pairs at \GTFullPrecision{}
precision. The \emph{class} level, the per-class rule the call graph replaced, asserts \GTClassExcess{} pairs
outside it at one site.
The call graph was written after that finding and is scored on the same ground truth.
The clean figure is therefore consistent with the repair but does not test it independently. \item The end-to-end check asks the same half of the question of the navigations two applications'
own suites perform. On realworld it does not discriminate.
The shell-rule level, which stands for earlier work, already contains all \EtoERealworldPairs{}
observed pairs (\EtoERealworldHitRooted{}).
On mifos the same level holds \EtoEMifosHitRooted{} of \EtoEMifosPairs{}, as
many as the full analysis (Table~\ref{tab:earlier}).
The hits and misses are in Section~\ref{sec:res-e2e}.
\item The rendered-interface check covers \ScopeApps{} of the six applications and answers half of the
question.
It ran on an earlier graph than the final call graph.
It ran on \ScopeJudged{} of \ScopeScreens{} screens (Section~\ref{sec:res-scope}). \end{itemize}
What remains unvalidated is the other half and the reach.
A site attributed to too \emph{many} screens produces edges whose witnesses all pass, because the
construct really is there and really names that target.
The rendered-interface check cannot convict it, since a predicted link missing from the DOM is equally a link rendered
only to a signed-in user.
The direction is unrefuted but close to untestable by this method.

\textbf{Construct validity.}
The denominator of must conversion (Section~\ref{sec:metrics}) over-counts, since it contains constructs in code the application never mounts.
It also under-counts, as follows.
The figure is therefore not a bound in either direction.
The independent route parser recognises navigation by four syntactic patterns: a \texttt{routerLink} attribute or
binding, and a call to \texttt{navigate} or \texttt{navigateByUrl} on a receiver whose name contains
\emph{router}.
Navigation expressed any other way is counted as neither a loss nor a success.
This includes helper-wrapped navigation, assignments to \texttt{location.href}, server-driven redirects and navigation inside third-party components.
The ratio is therefore recall against an enumerable syntactic proxy for an application's navigation,
and the completeness of that proxy is not itself measured.
The test-source cross-check is the only method that sees past it, because a URL the application's own
end-to-end specifications name is an address its authors expect to exist.
It reads the specifications without running them, so a URL there is named and not observed, and it
covers one application.

Node coverage carries a stronger form of the threat, in the opposite direction.
Its denominator is the screens the independent route parser can address, so a screen whose path neither tool can read
is outside the ratio and not counted against it.
Refusing to read more raises reported coverage.

The witness check judges every location an edge cites.
A construct therefore counts towards must conversion only on evidence the check has read.

The loss breakdown has a second construct threat: the precedence order of the rules is part of the
measure.
Testing the target first would make the target category larger than attribution on
\AttrOrderFlips{} applications (\AttrOrderFlipWhere{}).
Category sizes are comparable within one application and under one order, which is what
Chapter~\ref{ch:results} claims.

\textbf{Internal validity.}
Extraction runs on source trees whose dependencies are not installed, so the TypeScript compiler cannot
resolve types from the framework or third-party packages.
Type-directed steps such as narrowing an injected receiver are weaker, and the figures are
properties of that configuration.
The mutation experiment has small samples per operator, in single digits for several operator and
application pairs.
They were drawn with a seeded quota, so the pattern of which operators are detected is more
trustworthy than the rates \cite{papadakis2019mutation}.
The exclusion rule uses the graph itself, which couples the sample to the tool's own recovery.
The excluded mutants are reported, but the sample is not independent of the artefact under test.
The test-source cross-check measures node existence and nothing finer.

\textbf{External validity.}
The corpus is six open-source Angular applications pinned to single commits.
Proprietary applications may differ in size, routing style and navigation mechanisms
\cite{kalliamvakou2014promises,kalliamvakou2016indepth}.
The selection is tool-driven.
\texttt{SAP/spartacus} was rejected because its routing is configuration-driven in a way the
extractor does not support (Section~\ref{sec:corpus}).
Figures measured on the corpus are therefore likely optimistic for Angular applications at large.
The evidence from outside the extractor's own reading is uneven.
URLs named in end-to-end test sources (read and not run) come from alfresco alone.
The runtime router configuration was read on all \ProbeApps{} applications.
Rendered links were observed, on an earlier graph and in part, on \ScopeApps{}.
The end-to-end suites were observed on realworld and mifos, covering
\EtoEMifosPairs{} distinct transitions of an application with \RoundTwoRoutesMifos{} route entries.
The findings about scale therefore rest mainly on the runtime configuration and the extractor's own
reading.
The results are Angular results.
The graph model and the lattice do not mention Angular, and Section~\ref{sec:transfer} argues about
other routers.
That argument, however, is not a measurement.

\textbf{Oracle independence.}
The independent route parser shares no code with the extractor, which makes it independent and says nothing about
whether it is right.
The independence is also partial (Table~\ref{tab:repairs}).
For the route-path forms added during development (a table behind a freezing call, a path assembled
from constants, a module reached through a workspace alias), the two tools apply the same rules.
Their agreeing about such a path is weaker evidence than their agreeing about a path each reads its own
way.
A differential test only shows that two implementations disagree, never which is right.
Each disagreement was resolved by reading the source, by the person who wrote both implementations.

\textbf{Experimenter bias.}
I wrote the extractor, designed the evaluation, selected the corpus and defined the defect classes.
The mitigations applied are these.
\begin{itemize}[nosep]
\item The six selection criteria were fixed before any results, and rejected candidates are reported
with their reasons.
One rejection was on tool support, a ground outside the criteria.
Two retained applications do not meet the test-suite criterion (Section~\ref{sec:corpus}).
\item The witness check is mechanical and re-parses the source at the claimed location. It checks the
witness and not the address.
An edge whose target node sits at a path the application does not mount satisfies it.
\item Commits and harness version are pinned, and the result tables are generated from result files (Section~\ref{sec:protocol}).
The two tables compiled by hand are the defect-class table and the table of the witness check and the
test-source cross-check.
\item The protocol corrections are reported in full, including those that made the figures worse
(Appendix~\ref{app:repairs}).
\end{itemize}
Several controls were not applied:
\begin{itemize}[nosep]
\item no blinding, no second annotator for the loss classification, and no independent audit of the
residual constructs inspected by hand;
\item the defect classes were defined after the first observations, not in advance;
\item the source-level reading of the final must tier (its \CensusMustSites{} constructs) was made by the artificial intelligence (AI)
assistant declared at the front, which also helped write the extractor and the independent route parser.
It was made with the extractor's decisions in view and no human second reader checked it.
Its agreement with the extractor therefore shows consistency with the stated rules, not independent correctness;
\item uigraph was developed on, and is measured on, the same six
applications, so the figures are fitted to this corpus and are not estimates for other applications.
\end{itemize}
The requirement of two lines of evidence per defect class is the strongest control on the last point,
and the reason one class is still a candidate.
It is weaker than it appears.
Table~\ref{tab:defects} shows that four of the five accepted classes rest on one named method, the
rest of their evidence being observations.
For those four the control is therefore one method plus my own judgements.

\textbf{The runtime route oracle.}
The runtime route oracle of Section~\ref{sec:res-probe} removes the threat that mattered most, that two readers
of the same route files agree on a wrong address and the agreement is read as a result.
It brings in four others.
\begin{enumerate}[nosep]
\item A selection effect: the method needs an application that builds and boots, and a corpus chosen
partly for being buildable is not chosen at random.
\item It reads the route configuration and not the rendered page, so a route behind a guard that always
refuses, or whose component throws on render, counts as mounted.
\item The build was a development build in each case, because Angular removes the debug surface the
method reads in production mode.
A production configuration can differ in ways a route table does not show.
\item Where a workspace application composes the library a measurement is scoped to, the framework's
set is a superset of the library's, and part of any difference is scope and not reading. Two of
alfresco's three differences are of this kind.
\end{enumerate}

These threats do not make a mounted address look unmounted; they weaken the claim in the other direction.
The claim is that an address the extractor mints is one the framework mounts, tested by finding the
extractor's set inside the framework's. Threats 2 and 4 make the framework's set \emph{larger} than what
a visitor can reach, which makes containment easier to satisfy. Threat 3 moves the set in no fixed
direction, since a production configuration can add routes as well as drop them.
Both results are statements about the development router configuration that was recorded, under the
abstraction of Section~\ref{sec:res-probe}.
The compared graph mints \ProbeInvented{} such addresses, which is the absence of a failure on that
configuration and not a proof of its impossibility.

\section{Declared approximations}

The constructs the analysis does not decide are listed in Appendix~\ref{sec:approximations}. The list
separates constructs read exactly, where a failure is a defect Chapter~\ref{ch:results} counts, from
constructs approximated, where no completeness is claimed. A declared approximation can still cost
transitions the suites take; the end-to-end check measures that cost on mifos
(Section~\ref{sec:res-e2e}).

\section{Transfer to other frameworks}
\label{sec:transfer}

All measurements in this thesis concern Angular.
This section argues how much of the approach would carry over to React with React Router and to Vue
with Vue Router \cite{reactrouterdocs,vuerouterdocs}.
It is a qualitative analysis of the routers' documented behaviour, not a measurement; showing that the approach performs as well on React or Vue would require a corpus for each framework.

The analysis has a framework-independent part and a framework-specific part. The
graph model, the lattice of scopes, the fixpoint, anchored edges and the expansion
query do not mention Angular. What is specific to Angular are the four input relations
of Section~\ref{sec:attribution}: component nesting, router outlets, the call graph
with dependency injection, and dispatch relations. Moving to another framework means
deriving the same four relations from that framework's constructs.

Outlets transfer directly: in all three routers a parent route renders a
placeholder where the matched child is inserted. The components above the placeholder
form a shell that is present on every child screen, so the scopes $\top$ and
$\mathit{subtree}(r)$ have the same meaning.

The call graph is easier in React and Vue, because a hook or composable gives every navigation site an owning component and no injector needs to be resolved. Code that imports the router instance directly is the analogue of Angular's service case and can be resolved by ordinary import resolution, as in the JavaScript call graphs of Feldthaus et al.\ \cite{feldthaus2013jscg}.

Redirects do not transfer, and the asymmetry was found by porting the model's reading. Angular's validator refuses a route that declares
\texttt{redirectTo} beside \texttt{children} or a \texttt{component}.
A redirecting Angular route therefore has neither a layout nor a component to lose, and the model's
step from ``redirects'' to ``renders nothing'' costs nothing. Vue Router applies a record's redirect only when that record
is the last one matched, so a redirecting parent with children does render its component for a
matching child. Carrying the Angular reading into the Vue extractor dropped those components.
The extractor now keeps the component of a redirecting record that has children. It declares
one residue as a note at the parent rather than modelling it as an edge: navigating to the
parent's own address lands on the redirect target with that component unmounted. The
same porting exposed a second Vue-only hazard.
A table that mixes the Vue~2 catch-all \texttt{'*'} with the Vue~3 form \texttt{'/:pathMatch(.*)*'}
folds both to one address.
The extractor keeps the record declared first. Where that is the Vue~2 form, which Vue Router~4 rejects, the surviving component
is the wrong one. The collision is declared, but which record to keep is not something the table alone
determines.

Guards move to different scopes. A React Router loader redirect plays the role of \texttt{canActivate}, so its scope is the join of the origins of the navigations that may run the loader, as the relation $T$ gives for entry guards, and not $\mathit{subtree}(r)$. Vue's global \texttt{beforeEach} guard corresponds most closely to a \texttt{canActivateChild} on a root route with no component, anchored at $\top$.

Component nesting is harder in React. Angular and Vue keep templates in a
separate language, so the nesting relation can be read from element names. In React,
components are created in JSX, which is ordinary JavaScript. A component can be stored in a variable, passed as a property or as \texttt{children}, and rendered
under any condition the language can express. Recovering the nesting relation then becomes a data flow
problem closer to points-to analysis than to reading a template.
This is the main place where a React version would need new work.

The lattice, the fixpoint and the anchored representation transfer unchanged.
Of the relations, outlets and guards transfer with small changes, the call graph becomes simpler,
and component nesting becomes harder in React.
These are expectations derived from the routers' documented behaviour, and testing them on real
applications is left for future work (Section~\ref{sec:future}).

\section{Future work}
\label{sec:future}

This section looks forward; the development versions and the corrections between them are listed in
Appendix~\ref{app:repairs}.
The topics are listed in the order of how much each would change what the thesis can claim.

\paragraph{An oracle for the screen set on a large application.}
This is the cheapest addition that would strengthen the thesis most.
The hand-built ground truth of Section~\ref{sec:res-groundtruth} covers only the smallest
application.
Reading every site by hand is not practical on an application of hundreds of screens,
where most of the corpus's (site, screen) attributions are made (Table~\ref{tab:scopebaseline}).
Two routes lead there: driving such an application into a signed-in state and reading its
rendered links, or annotating a stratified sample of its sites instead of all of them.

\paragraph{The unimplemented part of the model.}
Chapter~\ref{ch:analysis} specifies two things the prototype leaves out.
\begin{itemize}[nosep]
\item The dispatch relations that carry NgRx effects and router event subscriptions.
\item The injector hierarchy that the implemented call graph reads past, since it takes every
provider binding as applying application-wide (Section~\ref{sec:approximations}). Reading
bindings relative to the injector that resolves them would narrow the scopes, and the fixpoint
already covers it. The call graph follows class hierarchy analysis in the style the
framework-modelling literature recommends \cite{dean1995cha,sridharan2011f4f,antoniadis2020enterprise}.
\end{itemize}
Guards that return a \texttt{UrlTree} belong to the same group, since a redirecting guard is itself
a navigation site and not only a label.
The wider treatment of guards in Chapter~\ref{ch:analysis}, where the guard of an edge is the
conjunction of the conditions on the containment path with the guards the router runs, is the more
substantial piece of work.
Its test is the guard-removal operator, which on this corpus reaches only
guards the graph already records on an edge (Section~\ref{sec:res-mutation}).

\paragraph{A runtime tier.}
The runtime tier is the complement to the whole static approach and not an extension of it.
The model already admits observation identifiers as witnesses beside source locations.
Setting that tier against the dynamic lineage, where the state flow graph comes from firing events
on a running application \cite{crawljax2012,feedex2013,fraggen2022}, would answer the question
this thesis cannot.
The question is how much of what a crawler finds the static graph already contains, and how
much of what the static graph contains a crawler never reaches.
The Android work combines static analysis with runtime traces in the same way \cite{promal2022,goalexplorer2019}.

For applications used signed in, such a combination is also the way to complete the graph: a signed-in run against a seeded backend would show which may edges occur, and an observed transition would gain an observation witness beside its source witness.
That is evidence that the may edge exists, not a promotion to must, since a must edge asserts the
transition on every execution the claim ranges over and one run shows it on one.
A candidate never observed stays may.
The combination can therefore add evidence without weakening Claim~\ref{claim:must}.
This thesis neither builds nor measures it.

\paragraph{A larger and broader corpus.}
An automated search over a larger pool of repositories, keeping applications the extractor handles badly \cite{kalliamvakou2014promises,kalliamvakou2016indepth}, would weaken the tool-driven selection bias of Section~\ref{sec:limitations}.
A case study on a proprietary industrial application, reporting only aggregate figures, is the most direct answer to the external validity threat.

\paragraph{The use that motivated the thesis.}
That use is still unmeasured.
Consider a graph of screens and guarded transitions, recomputed from source and carrying a witness
on every edge.
Whether such a graph improves mobile testing agents or coding agents (Chapter~\ref{ch:introduction}), or improves automated test generation, is a
controlled experiment that this thesis does not attempt.
The graph and the harness are the artefacts such an experiment would need, and the defects quantified in
Chapter~\ref{ch:results} are what it would have to control for.

\section{Summary}

The transfer argument is a hypothesis, not a measurement.
The most valuable next steps are an oracle for the screen set on a large application and a runtime tier
that completes the graph for signed-in use.
